\section{Билет 53. Распределение Максвелла молекул по величине скорости. График функции распределения Максвелла. Изменение формы графика при изменении температуры и массы молекул с обоснованием.}

\begin{definition}
\textbf{Распределение Максвелла} описывает вероятность того, что молекула идеального газа, находящегося в термодинамическом равновесии при температуре $T$, имеет модуль скорости в интервале $(v, v + dv)$.
\par\noindent Функция распределения по модулю скорости имеет вид:
\begin{equation}
    F(v) = 4\pi \left( \frac{m}{2\pi k T} \right)^{3/2} v^2 \exp\left( -\frac{m v^2}{2 k T} \right), \label{eq:maxwell_speed_53}
\end{equation}
где:
\begin{itemize}
    \item $m$ --- масса одной молекулы;
    \item $k = 1.38 \cdot 10^{-23}$ Дж/К --- постоянная Больцмана;
    \item $T$ --- абсолютная температура.
\end{itemize}
Вероятность того, что модуль скорости молекулы находится в интервале $(v, v + dv)$, равна:
\begin{equation}
    dP = F(v) \, dv.
\end{equation}
\end{definition}

\begin{property}[Структура функции $F(v)$]
Функция \eqref{eq:maxwell_speed_53} представляет собой произведение двух множителей:
\begin{enumerate}
    \item $v^2$ --- геометрический множитель, возникающий при переходе от декартовых координат $(v_x, v_y, v_z)$ к сферическим в пространстве скоростей. Объём сферического слоя радиуса $v$ и толщиной $dv$ равен $4\pi v^2 dv$.
    \item $\exp(-mv^2/2kT)$ --- больцмановский множитель, определяющий вероятность того, что молекула имеет кинетическую энергию $E = mv^2/2$.
\end{enumerate}
При $v \to 0$ функция $F(v) \to 0$ (из-за множителя $v^2$), при $v \to \infty$ функция экспоненциально убывает. Следовательно, $F(v)$ имеет максимум при некоторой конечной скорости.
\end{property}

\begin{remark}
\textbf{График функции распределения $F(v)$.}
\begin{itemize}
    \item Кривая начинается в нуле, возрастает, достигает максимума при наиболее вероятной скорости $v_{\text{вер}}$, затем монотонно убывает, асимптотически приближаясь к нулю.
    \item Площадь под кривой равна единице (условие нормировки): $\int_0^\infty F(v) \, dv = 1$.
    \item Физический смысл $F(v)$: доля молекул, модули скоростей которых лежат в единичном интервале вблизи значения $v$.
\end{itemize}
\end{remark}

\begin{property}[Влияние температуры на форму графика]
При увеличении температуры $T$ (при фиксированной массе $m$):
\begin{enumerate}
    \item Максимум кривой смещается вправо: $v_{\text{вер}} = \sqrt{2kT/m} \propto \sqrt{T}$.
    \item Высота максимума уменьшается: $F(v_{\text{вер}}) \propto 1/\sqrt{T}$.
    \item Кривая становится шире и более ``пологой'': разброс скоростей увеличивается.
\end{enumerate}
\textbf{Обоснование:} При повышении температуры средняя кинетическая энергия молекул $\langle E \rangle = 3kT/2$ растёт. Молекулы в среднем движутся быстрее, поэтому распределение ``размазывается'' в область больших скоростей. Условие нормировки требует, чтобы площадь под кривой оставалась равной 1, поэтому при увеличении ширины высота максимума должна уменьшаться.
\end{property}

\begin{property}[Влияние массы молекул на форму графика]
При увеличении массы молекулы $m$ (при фиксированной температуре $T$):
\begin{enumerate}
    \item Максимум кривой смещается влево: $v_{\text{вер}} \propto 1/\sqrt{m}$.
    \item Высота максимума увеличивается: $F(v_{\text{вер}}) \propto \sqrt{m}$.
    \item Кривая становится уже и более ``острой'': разброс скоростей уменьшается.
\end{enumerate}
\textbf{Обоснование:} При одинаковой температуре средняя кинетическая энергия всех молекул одинакова: $\langle mv^2/2 \rangle = 3kT/2$. Следовательно, более тяжёлые молекулы должны иметь меньшие скорости, чтобы их кинетическая энергия оставалась той же. Поэтому распределение для тяжёлых молекул сосредоточено в области меньших скоростей.
\end{property}

\begin{remark}
\textbf{Универсальная форма распределения.} Если ввести безразмерную скорость $u = v / v_{\text{вер}}$, то функция распределения принимает вид:
\begin{equation}
    f(u) = \frac{4}{\sqrt{\pi}} u^2 e^{-u^2},
\end{equation}
который не зависит ни от рода газа, ни от температуры. Это означает, что все кривые Максвелла подобны друг другу и отличаются только масштабом по осям.
\end{remark}

\begin{figure}[htbp]
    \centering
    \begin{tikzpicture}[
        scale=1.1,
        >=Stealth,
        font=\small,
        thick
    ]

    \node[font=\bfseries] at (4, 4.6) {Распределение Максвелла по модулю скорости $F(v)$};

    % === ОСИ ===
    \draw[->] (0, 0) -- (7, 0) node[right, font=\normalsize] {$v$};
    \draw[->] (0, 0) -- (0, 4) node[above, font=\normalsize] {$F(v)$};
    \node[below left] at (0, 0) {$O$};

    % === КРИВЫЕ ДЛЯ РАЗНЫХ ТЕМПЕРАТУР ===
    % F(v) = C * v^2 * exp(-a*v^2)
    % Используем \x*\x вместо \x^2 для корректного парсинга

    % T1 (низкая температура) — узкий высокий пик
    \draw[blue, very thick, domain=0:6, samples=200]
        plot (\x, {4.9*\x*\x*exp(-0.6*\x*\x)});

    % T2 (средняя температура) — средний пик
    \draw[green!60!black, very thick, domain=0:6, samples=200]
        plot (\x, {2.2*\x*\x*exp(-0.35*\x*\x)});

    % T3 (высокая температура) — широкий низкий пик
    \draw[red, very thick, domain=0:6, samples=200]
        plot (\x, {0.92*\x*\x*exp(-0.2*\x*\x)});

    % === МАКСИМУМЫ КРИВЫХ ===
    % T1: максимум при v ≈ 1.29
    \fill[blue] (1.29, 3) circle (3pt);
    \draw[dashed, blue] (1.29, 0) -- (1.29, 3);
    \node[blue, below, font=\footnotesize] at (1.29, -0.2) {$v_{\text{1}}$};

    % T2: максимум при v ≈ 1.69
    \fill[green!60!black] (1.69, 2.3) circle (3pt);
    \draw[dashed, green!60!black] (1.69, 0) -- (1.69, 2.3);
    \node[green!60!black, below, font=\footnotesize] at (1.69, -0.2) {$v_{\text{2}}$};

    % T3: максимум при v ≈ 2.24
    \fill[red] (2.24, 1.7) circle (3pt);
    \draw[dashed, red] (2.24, 0) -- (2.24, 1.7);
    \node[red, below, font=\footnotesize] at (2.24, -0.2) {$v_{\text{3}}$};

    % === ПОДПИСИ КРИВЫХ ===
    \node[blue, font=\footnotesize] at (2.8, 0.9) {$T_1$};
    \node[green!60!black, font=\footnotesize] at (3.5, 0.9) {$T_2$};
    \node[red, font=\footnotesize] at (4.35, 0.9) {$T_3$};

    % === СТРЕЛКИ ПОКАЗА НАПРАВЛЕНИЯ ИЗМЕНЕНИЯ T ===
    \draw[->, purple, thick] (1.5, 3.5) -- (2.2, 3.5);
    \node[purple, above, font=\footnotesize] at (1.85, 3.5) {$v_{\text{вер}}$ растёт};

    \draw[->, purple, thick] (4.5, 2.5) -- (4.5, 1.5);
    \node[purple, right, font=\footnotesize] at (4.5, 2.0) {$F_{\text{max}}$ убывает};

    % === БЛОК СВОЙСТВ ===
    \node[draw, fill=yellow!10, rounded corners, align=left, font=\small, inner sep=5pt] at (11, 3.2)
        {Свойства $F(v)$:\\[4pt]
        $\bullet$ $F(0) = 0$\\[4pt]
        $\bullet$ $F(v) \to 0$ при $v \to \infty$\\[4pt]
        $\bullet$ $v_{\text{вер}} = \sqrt{\frac{2kT}{m}} \propto \sqrt{T}$\\[4pt]
        $\bullet$ $F_{\text{max}} \propto \frac{1}{\sqrt{T}}$\\[4pt]
        $\bullet$ $\int_0^\infty F(v) dv = 1$};

    \end{tikzpicture}
    \caption{Распределение Максвелла по модулю скорости $F(v)$ для трёх температур $T_1 < T_2 < T_3$. При увеличении температуры максимум смещается вправо ($v_{\text{вер}} \propto \sqrt{T}$) и становится ниже ($F_{\text{max}} \propto 1/\sqrt{T}$). Кривая расширяется, показывая больший разброс скоростей. Площадь под каждой кривой равна 1 (условие нормировки).}
    \label{fig:maxwell_speed_distribution_fixed}
\end{figure}

